\section{Recovering the modulus and twist from the signed zero spectrum}
\label{sec:packet}
The residues in \texttt{New.isPolar\_primeRemainder} record only the
difference between two sums of zero multiplicities. Contributions at a common
zero can cancel exactly, and the names of the characters have disappeared.
Nevertheless the signed data determine both arithmetic parameters. The proof
also identifies the limiting oscillations of the prime remainder used in the
Lean argument.

Fix $8\mid q$, $q>0$, and a prime $d\equiv7\pmod 8$ coprime to $q$.
Write $h=\varphi(q)$ and $\eta(n)=(n/d)$. For each $\chi\bmod q$ let
$\chi^*$ be its inducing primitive character. For a real ordinate $\gamma$ put
$\rho_\gamma=1/2+i\gamma$ and set
\[
 m(\gamma)=\frac1h\sum_{\chi\bmod q}
 \left(\operatorname{ord}_{\rho_\gamma} L(s,\chi^*)-
       \operatorname{ord}_{\rho_\gamma} L(s,\chi^*\eta)\right),
 \qquad b_\gamma=\frac{m(\gamma)}{1/4+\gamma^2}.
\]
Removing imprimitive Euler factors preserves orders in the open critical
strip. Thus $b_\gamma$ is precisely the residue in the source. Dirichlet GRH
puts every nontrivial zero in these sums on the specified line.
The input to the inverse problem is the signed coefficient at each frequency,
with coincident frequencies combined.

\paragraph{Recovery theorem.}
The series $F(t)=\sum_\gamma b_\gamma e^{i\gamma t}$ is real, even and
uniformly convergent. Its derivative has a jump $\log d$ at zero. At $\log p$
for every prime $p\nmid qd$ it has a jump precisely if
$p\equiv1\pmod q$ and $\eta(p)=-1$, in which case the jump is
$2\log p/\sqrt p$. Consequently the data $\{(\gamma,b_\gamma)\}$
determine $(q,d)$ even if they have been changed at finitely many frequencies.
More generally, the same reconstruction works after adding any continuously
differentiable function to $F$. It concerns the full infinite spectrum and
does not assert a cutoff at which a finite list determines the parameters.

\paragraph{Proof of the jump formula.}
The classical zero count $N(T)=O_{q,d}(T\log(2+T))$ ensures
$\sum_\gamma |b_\gamma|<\infty$ by partial summation. Both families are closed
under complex conjugation, whence $b_{-\gamma}=b_\gamma\in\mathbb R$.
For a prime power $n=p^k$ let $Q_p=p^{v_p(q)}$, $q^{(p)}=q/Q_p$, and put
\[
 a(p^k)=\frac{1-\eta(p^k)}{\varphi(Q_p)}
               \mathbf 1_{p^k\equiv1\bmod q^{(p)}} .
\]
In distributions on the real line the following identity holds, with
$\delta_u$ a unit point mass at $u$:
\begin{align}
 F''-\tfrac14 F={}&(\log d)\delta_0-
       \frac2h\cosh(t/2)\,dt \notag\\
 &+\sum_{p^k\ge2}\frac{\log p}{p^{k/2}}\,a(p^k)
                     (\delta_{\log p^k}+\delta_{-\log p^k}).
 \label{eq:packet-ode}
\end{align}
Here derivatives are taken in the distributional sense and functions are
identified with their densities. To check the constants, test against
$g\in C_c^\infty(\mathbb R)$ with
$H(u)=\int g(t)e^{iut}\,dt$. The Guinand--Weil explicit formula for a
primitive character of conductor $f$ contributes $\log(f/\pi)g(0)$, the
gamma term for its parity, and a negative sum of prime powers with
coefficients $\Lambda(n)\chi^*(n)/\sqrt n$ at one sign of $\log n$ and
the conjugate coefficients at the other. The primitive principal character
also contributes $H(i/2)+H(-i/2)$ from the poles at $0,1$ of the completed
zeta function. This is the normalization of
A.~Chirre, V.~J.~Pereira J\'unior and D.~de Laat,
\emph{Primes in arithmetic progressions and semidefinite programming},
Math. Comp. 90 (2021), \S2.1,
\href{https://doi.org/10.1090/mcom/3638}{doi:10.1090/mcom/3638}.
Equivalently the formula follows by shifting the contour of the logarithmic
derivative of the completed primitive $L$-function.

As $q>2$, half the characters have each parity. The odd character $\eta$
interchanges the parities. All gamma terms therefore cancel between the two
families. The two conductors for $\chi^*$ and $\chi^*\eta$ are $f$ and
$fd$, since $(f,d)=1$. Only the untwisted family contains the primitive
principal character. At $p^k$ the average of the untwisted primitive
character values is $\mathbf 1_{p^k\equiv1\bmod q^{(p)}}/\varphi(Q_p)$.
Indeed a nonprincipal $p$-component makes the primitive value zero,
and orthogonality over the characters modulo $q^{(p)}$ evaluates the
remaining sum. Multiplication by $1-\eta(p^k)$ gives the displayed $a$,
also at $p=d$. Finally $F''-F/4=-\sum_\gamma m(\gamma)e^{i\gamma t}$,
which accounts for the signs in \eqref{eq:packet-ode}.

The right side is smooth apart from a locally finite set of point masses.
Integrating across one such point shows that the jump of $F'$ equals its mass.
Away from the atoms $F$ satisfies an ordinary differential equation with
analytic coefficients, hence is analytic there. Evenness then gives
$F(t)=F(0)+\tfrac12\log d\,|t|+O_{q,d}(t^2)$ near zero.

\paragraph{Recovering $q$.}
Let $S$ contain the primes $p$ at which $F'$ has a positive jump at $\log p$.
For all sufficiently large $X$,
\[
 q=\gcd\{p-1:p\in S,\ p>X\}.
\]
Every entry is divisible by $q$ after the finitely many primes dividing $qd$
have been discarded. Suppose a prime $\ell$ divided the gcd after division by
$q$. We produce infinitely many primes in $S$ with $\ell\nmid(p-1)/q$.
If $\ell\mid q$, require $p\equiv1+q\pmod{q\ell}$ and independently choose a
nonzero quadratic nonresidue modulo $d$. If $\ell\nmid qd$, then $\ell$ is odd.
Choose a unit other than $1$ modulo $\ell$, the residue $1$ modulo $q$, and a
nonresidue modulo $d$. If $\ell=d$, the conditions $p\equiv1\pmod q$ and
$\eta(p)=-1$ already imply $p\not\equiv1\pmod\ell$. Each construction gives
a reduced class by the Chinese remainder theorem and infinitely many primes
by Dirichlet's theorem. Thus no $\ell$ divides the quotient. The jump at zero
recovers $d$ by exponentiation. All these jumps survive a $C^1$ perturbation,
including any finite trigonometric polynomial.\hfill$\square$

\paragraph{Identifying the spectral function on actual prime sums.}
Define the remainder exactly as in \texttt{New.primeRemainder},
\[
 R(x)=
 \sum_{\substack{n\le x\\n\equiv1\bmod q}}
 \Lambda(n)(1-\eta(n))(1/n-1/x)-h^{-1}\log x-c,
\]
where $c$ gives $R(x)\to0$ by the prime number theorem for fixed progressions.
The derivative jumps determine
\begin{equation}
 \sqrt{x}\,R(x)=F(\log x)+O_{q,d}\!\left(\frac{1+\log x}{\sqrt{x}}\right)
 \qquad(x\ge 1).
 \label{eq:packet-remainder}
\end{equation}
In particular this normalization remains bounded. To prove the identity, use
$F'(0+)=(\log d)/2$ and solve \eqref{eq:packet-ode} on $t\ge0$. Write
$S^*(x)=\sum_{p^k\le x}\log(p)\,a(p^k)(p^{-k}-x^{-1})$. The homogeneous
solutions are $\cosh(t/2)$ and $\sinh(t/2)$, the particular solution for
the continuous forcing is $-2t\sinh(t/2)/h$, and an atom of mass $B$
at $u>0$ adds $2B\sinh((t-u)/2)$ for $t\ge u$. Consequently
\[
 \frac{F(\log x)}{\sqrt{x}}=
 S^*(x)-\frac{1-x^{-1}}h\log x+
 \frac{F(0)}2(1+x^{-1})+\frac{\log d}2(1-x^{-1}).
\]
The difference between $S^*(x)$ and the sum in $R(x)$ comes exclusively from
powers of $p\mid q$. It is
$C_{\rm bad}+O_q((1+\log x)/x)$, where
$C_{\rm bad}=\sum_{p\mid q}\sum_{k\ge1}\log(p)\,a(p^k)p^{-k}$.
This follows by bounding a geometric tail for each such $p$ and at most
$\log x/\log p$ terms of size $2\log(p)/x$. Taking $x\to\infty$ gives
$c+C_{\rm bad}+(F(0)+\log d)/2=0$, proving \eqref{eq:packet-remainder}.

For the parameters used in the lattice proof, $18480\mid q$ and
$d\equiv q/2-1\pmod q$, the source proves $\eta(p)=1$ for every $p\mid q$
(\texttt{New.Parameters.chiK\_eq\_one\_of\_dvd}). Then $C_{\rm bad}=0$
and every correction at $p\mid q$ vanishes. In this case the comparison
is exact for every $x\ge1$:
\[
 F(0)=-2c-\log d,\qquad
 \sqrt{x}R(x)=F(\log x)+\frac{c+\log d-h^{-1}\log x}{\sqrt{x}} .
\]
For example, $(q,d)=(18480,9239)$ satisfies the source's parameter conditions.
The two prime jumps at $258721$ and $794641$ have
$p-1=14q$ and $p-1=43q$, respectively. Their gcd is already $q$.
Trial division and Euler's criterion verify both primes and both nonresidues.

\paragraph{Oscillation at the critical scale.}
The logarithmic mean of $\sqrt{x}\,R(x)$ is $b_0$ by termwise averaging.
The cusp at zero implies that infinitely many nonzero frequencies survive.
For every $\gamma\ne0$ with $b_\gamma\ne0$,
\[
 \liminf_{x\to\infty}\sqrt{x}\,R(x)\le b_0-|b_\gamma|,
 \qquad
 \limsup_{x\to\infty}\sqrt{x}\,R(x)\ge b_0+|b_\gamma|.
\]
Indeed the averages of $F(t)$ times $1\pm\cos(\gamma t)$ over $[0,T]$
tend to $b_0\pm b_\gamma$. Each weight is nonnegative and has mean one.
Equation~\eqref{eq:packet-remainder} transfers the bounds to the prime sum.
The mean includes any genuine central zero. This argument uses neither
simplicity nor linear independence of zero ordinates.

\paragraph{A finite-height approximation with an error rate.}
With coefficients known through height $T$, write $F_T(t)=\sum_{|\gamma|\le
T}b_\gamma e^{i\gamma t}$. For any fixed point $u$, let
$J(u)=F'(u+)-F'(u-)$, taking $J(u)=0$ where $F$ is smooth. Define
\[
 \widehat J_T(u)=
 \frac{F_T(u+z)-2F_T(u)+F_T(u-z)}{z},\qquad
 z=\sqrt{\log(2+T)/T}.
\]
Then $\widehat J_T(u)=J(u)+O_{q,d,u}(\sqrt{\log(2+T)/T})$ as
$T\to\infty$. In fact the zero count gives
$\lVert F-F_T\rVert_\infty=O_{q,d}(\log(2+T)/T)$. The full second
difference divided by $z$ equals $J(u)+O_{q,d,u}(z)$ by one-sided Taylor
expansion in an interval containing no other atom. Truncation costs at most
$4\lVert F-F_T\rVert_\infty/z$. At zero, rounding $\exp(\widehat J_T(0))$
eventually gives the exact integer $d$. For fixed primes away from $qd$
one may likewise distinguish the two possible jump sizes. Identifying a
uniform numerical stopping rule would require explicit bounds in these errors.

The inverse calculation uses the specific signed residue identity exposed by
the repository, classical explicit formulas, and character orthogonality.
It proves uniqueness for these aggregate data. The finite-height estimate is for exact ordinates and signed multiplicities.
The next estimate also permits quantified errors in the spectral data.
